% =====================================================================
% Mục 6 — Kết luận  (BẢN TIẾNG VIỆT)
%
% Bam bon phan duoc yeu cau: tom tat; mo hinh tot nhat cho tung kich ban
% su dung; y nghia bao ton; huong nghien cuu tiep theo.
%
% Theo scientific-writing.md va writing_quality_check.md cua skill
% academic-paper:
%   * Ket luan khong duoc gioi thieu du lieu hay lap luan moi chua xuat
%     hien o Muc 4 va Muc 5. Moi khang dinh o day deu tro ve mot tieu muc
%     cu the cua Muc 5.
%   * Khong dung dau gach ngang lam dau cau.
%   * Do dai doan phai bien thien, khong deu nhau.
%   * Cac khang dinh ve hien trang bao ton nam ngoai du lieu nghien cuu
%     nen duoc danh dau [CAN NGUON].
% =====================================================================

\section{Kết luận}
\label{vi:sec:conclusion}

Nghiên cứu này đánh giá đối sánh năm kiến trúc học sâu cho bài toán định danh
17 loài cây lá kim từ ảnh mặt cắt ngang ở thang đo vĩ mô, trên một phân hoạch
gom theo khóa nguồn và một registry đánh giá bất biến duy nhất. ConvNeXt-Tiny
đạt độ chính xác cao nhất và ResNet-50 kém rõ rệt, trong khi ba mô hình còn lại
chênh nhau những khoảng cách cùng bậc độ lớn với biến động giữa các seed, nên
thứ tự nội bộ của nhóm này chưa được xác lập
(Mục~\ref{vi:sec:discussion:ranking}).

Cấu trúc lỗi mang nhiều thông tin hơn bản thân con số độ chính xác. Toàn bộ khó
khăn còn lại của mô hình mạnh nhất nằm ở một cặp loài duy nhất trong cùng một
chi (Mục~\ref{vi:sec:discussion:cupressaceae}).

Lựa chọn mô hình phụ thuộc vào kịch bản sử dụng. Với hệ thống chạy trên máy
chủ, nơi chi phí tính toán không phải ràng buộc, ConvNeXt-Tiny là lựa chọn hợp
lý. Với định danh tại hiện trường trên điện thoại hoặc thiết bị nhúng,
MobileNetV3-Large phù hợp hơn, vì chỉ kém 1.19 điểm độ chính xác nhưng dùng
khoảng một phần sáu số tham số và một phần hai mươi khối lượng tính toán.
ResNet-50 không được khuyến nghị cho cả hai kịch bản, vì nó có chi phí tính
toán cùng bậc với ConvNeXt-Tiny mà độ chính xác thấp hơn bảy điểm.

Về ý nghĩa bảo tồn, bốn loài đặc hữu Việt Nam trong bộ dữ liệu đều được phân
loại tốt dưới mô hình mạnh nhất, trong đó hai loài đạt F1 tuyệt đối trên cả ba
seed (Mục~\ref{vi:sec:discussion:endemic})
[CẦN NGUỒN: hiện trạng bảo tồn của bốn loài này]. Một công cụ định danh tự động
đủ nhẹ để chạy tại hiện trường có thể hỗ trợ việc kiểm soát nguồn gốc gỗ và
giám sát khai thác. Kết luận đó chỉ có giá trị nếu hiệu năng được xác nhận lại
trên ảnh thu thập độc lập, điều mà nghiên cứu này chưa làm được
(Mục~\ref{vi:sec:discussion:external}).

Ba hướng tiếp theo là trực tiếp. Lặp lại mỗi ablation trên đúng ba seed như các
lần chạy chính sẽ làm các phép so sánh đó diễn giải được. Hoàn tất việc xác
minh khóa mẫu vật đối chiếu với sổ mẫu của bộ sưu tập sẽ chứng minh, thay vì
giả định, rằng phân hoạch tách được các mẫu vật vật lý. Và đánh giá trên ảnh
chụp bởi người khác, thiết bị khác và điều kiện ánh sáng khác là cách duy nhất
để tách hai cách đọc lành tính và bất lợi của kết quả thăm dò.
